%%%%%%%%%%%%%%%%%%%%%%%%%%%%%%%%%%%%%%%%%%%%%%%%%%%%%%%%%%%%%%%%%%%%%%%%%%%%%%%%%%
\begin{frame}[fragile]\frametitle{}
\begin{center}
{\Large Case Study: PCA on Car Specs}
\end{center}
\end{frame}

%%%%%%%%%%%%%%%%%%%%%%%%%%%%%%%%%%%%%%%%%%%%%%%%%%%%%%%%%%
\begin{frame}[fragile]\frametitle{The Data}
\begin{itemize}
\item 8 numeric specs for 426 cars from the SAS sample data set \texttt{SASHELP.CARS}: engine size, cylinders, horsepower, city and highway MPG, weight, wheelbase, length
\item Many features move together: Weight correlates 0.81 with engine size and $-0.79$ with highway MPG
\item Clean first, as in Session 16: 2 rotary-engine cars have a dot (\texttt{.}) for cylinders and are dropped
\item Here we use all 8 specs; Session 16 clustered on 5 of them
\end{itemize}
\begin{lstlisting}
import pandas as pd

cars = pd.read_csv('data/cars.csv', sep=';')
features = ['EngineSize', 'Cylinders', 'Horsepower', 'MPG_City',
            'MPG_Highway', 'Weight', 'Wheelbase', 'Length']
cars[features] = cars[features].apply(pd.to_numeric, errors='coerce')
cars = cars.dropna(subset=features)     # 428 -> 426 cars
X = cars[features]
\end{lstlisting}
\end{frame}

%%%%%%%%%%%%%%%%%%%%%%%%%%%%%%%%%%%%%%%%%%%%%%%%%%%%%%%%%%
\begin{frame}[fragile]\frametitle{Raw Features}
\begin{itemize}
\item PC1 explains 99.4\% of the variance
\item Its loading on Weight is 0.998; the other loadings are near 0
\item Weight is in the thousands, so its variance swamps the rest: PCA found the units, not the structure
\end{itemize}
\begin{lstlisting}
from sklearn.decomposition import PCA

pca = PCA()
pca.fit(X)

# In the Python shell:
>>> print(pca.explained_variance_ratio_[:3].round(3))
[0.994 0.005 0.   ]
>>> print(pca.components_[0].round(3))    # PC1 loadings, in feature order
[ 0.001  0.002  0.06  -0.005 -0.006  0.998  0.008  0.013]
\end{lstlisting}
\end{frame}

%%%%%%%%%%%%%%%%%%%%%%%%%%%%%%%%%%%%%%%%%%%%%%%%%%%%%%%%%%
\begin{frame}[fragile]\frametitle{Standardized Features}
\begin{itemize}
\item Standardize first (centering and scaling), then run PCA
\item PC1 explains 71.6\% of the variance, PC2 13.4\%
\item 3 components keep 92.4\%, 4 keep 95.1\%: 8 features become 3 or 4 numbers per car
\end{itemize}
\begin{lstlisting}
from sklearn.preprocessing import StandardScaler

Xs = StandardScaler().fit_transform(X)
pca = PCA()
pca.fit(Xs)

# In the Python shell:
>>> print(pca.explained_variance_ratio_.round(3))
[0.716 0.134 0.074 0.027 0.022 0.012 0.01  0.005]
>>> print(pca.explained_variance_ratio_.cumsum().round(3))
[0.716 0.849 0.924 0.951 0.973 0.985 0.995 1.   ]
\end{lstlisting}
\end{frame}

%%%%%%%%%%%%%%%%%%%%%%%%%%%%%%%%%%%%%%%%%%%%%%%%%%%%%%%%%%
\begin{frame}[fragile]\frametitle{Reading the Components}
\begin{itemize}
\item Loadings are the weights of each original feature in a component
\item PC1: engine size, cylinders, power and weight up, MPG down: overall size and power
\item PC2: wheelbase and length up, horsepower down: roominess versus power
\item The names are our reading; the overall sign of a component is arbitrary
\end{itemize}
\begin{lstlisting}
>>> load = pd.DataFrame(pca.components_[:2].T, index=features,
...                     columns=['PC1', 'PC2'])
>>> print(load.round(2))
              PC1   PC2
EngineSize   0.39 -0.09
Cylinders    0.37 -0.19
Horsepower   0.33 -0.40
MPG_City    -0.36  0.24
MPG_Highway -0.36  0.24
Weight       0.38  0.11
Wheelbase    0.32  0.57
Length       0.31  0.58
\end{lstlisting}
\end{frame}

%%%%%%%%%%%%%%%%%%%%%%%%%%%%%%%%%%%%%%%%%%%%%%%%%%%%%%%%%%
\begin{frame}[fragile]\frametitle{From Scratch with NumPy}
\begin{itemize}
\item The same steps as the by-hand example: covariance matrix, eigenvectors, sort by eigenvalue
\item \texttt{eigh} is made for symmetric matrices, such as a covariance matrix
\item The eigenvalues match scikit-learn's \texttt{explained\_variance\_}
\end{itemize}
\begin{lstlisting}
import numpy as np

C = np.cov(Xs.T)                     # 8 x 8 covariance matrix
eigvals, eigvecs = np.linalg.eigh(C)
order = np.argsort(eigvals)[::-1]    # largest eigenvalue first
eigvals, eigvecs = eigvals[order], eigvecs[:, order]

# In the Python shell:
>>> print(eigvals.round(3))
[5.738 1.073 0.597 0.218 0.177 0.092 0.082 0.042]
>>> print(np.allclose(eigvals, pca.explained_variance_))
True
\end{lstlisting}
\end{frame}

%%%%%%%%%%%%%%%%%%%%%%%%%%%%%%%%%%%%%%%%%%%%%%%%%%%
% Optional (Sep 2026 review): the SVD route is already taught in ml_pca.tex (Computing It with SVD).
% \begin{frame}[fragile]\frametitle{The Same Result via SVD}
% \begin{itemize}
% \item SVD works on the data directly: no covariance matrix is built
% \item The rows of \texttt{Vt} are the principal components; $\sigma^2/(n-1)$ are the eigenvalues
% \item The results match the eigenvalue route, up to the sign of each component
% \end{itemize}
% \begin{lstlisting}
% U, s, Vt = np.linalg.svd(Xs, full_matrices=False)

% # In the Python shell:
% >>> print((s ** 2 / (len(Xs) - 1)).round(3))
% [5.738 1.073 0.597 0.218 0.177 0.092 0.082 0.042]
% >>> print(np.allclose(np.abs(Vt[0]), np.abs(pca.components_[0])))
% True
% \end{lstlisting}
% \end{frame}

%%%%%%%%%%%%%%%%%%%%%%%%%%%%%%%%%%%%%%%%%%%%%%%%%%%
% Optional (Sep 2026 review): randomized solver, an advanced extra; restore for a scale-up discussion.
% \begin{frame}[fragile]\frametitle{Randomized PCA}
% \begin{itemize}
% \item Randomized PCA finds only the top $k$ components, using random projections
% \item Cost: $\mathcal{O}(n_{\max}^2 \cdot k)$ instead of $\mathcal{O}(n_{\max}^2 \cdot n_{\min})$ for the full SVD
% \item It is approximate, so set \texttt{random\_state}; use it when there are many samples or features and you need few components
% \end{itemize}
% \begin{lstlisting}
% pca_rand = PCA(n_components=3, svd_solver='randomized', random_state=0)
% pca_rand.fit(Xs)

% # In the Python shell:
% >>> print(pca_rand.explained_variance_ratio_.round(3))
% [0.716 0.134 0.074]
% \end{lstlisting}
% \end{frame}

%%%%%%%%%%%%%%%%%%%%%%%%%%%%%%%%%%%%%%%%%%%%%%%%%%%%%%%%%%
\begin{frame}[fragile]\frametitle{Project and Reconstruct}
\begin{itemize}
\item Project onto the first 2 components: 8 features become 2 numbers per car
\item Reconstruct: map the 2 numbers back to 8 features
\item The reconstruction error (0.151) equals the variance we dropped ($1 - 0.849$)
\end{itemize}
\begin{lstlisting}
W = eigvecs[:, :2]        # projection matrix: 8 x 2
Z = Xs @ W                # 2 numbers per car
Xr = Z @ W.T              # back to 8 features

# In the Python shell:
>>> print(Z.shape)
(426, 2)
>>> print(((Xs - Xr) ** 2).mean().round(3))
0.151
>>> print((1 - pca.explained_variance_ratio_[:2].sum()).round(3))
0.151
\end{lstlisting}
\end{frame}

%%%%%%%%%%%%%%%%%%%%%%%%%%%%%%%%%%%%%%%%%%%%%%%%%%%%%%%%%%
\begin{frame}[fragile]\frametitle{Quick Check: Reading Explained Variance}
\begin{block}{Think About It}
In the standardized cars data, PC1 explains 71.6\% of the variance. What does that mean?
\begin{enumerate}
  \item[A.] PC1 predicts a car's price with 71.6\% accuracy
  \item[B.] Keeping only PC1 retains 71.6\% of the total variance; the other 28.4\% is lost
  \item[C.] PC1 is the single most important of the original features
  \item[D.] 71.6\% of the cars are described perfectly by PC1
\end{enumerate}
\end{block}
\end{frame}

%%%%%%%%%%%%%%%%%%%%%%%%%%%%%%%%%%%%%%%%%%%%%%%%%%%%%%%%%%%%%%%%%%%%%%%%%%%%%%%%%%
\begin{frame}[fragile]\frametitle{Quick Check: Reading Explained Variance}
\begin{block}{Answer}
\textbf{Correct: B.} The explained variance ratio is the share of the total variance that a component carries. PC1 is a new axis, a mix of all eight features, not one of the original features (C). The ratio measures variance, not prediction accuracy (A) or the share of cars fitted (D).
\end{block}
\end{frame}
